\subsection{Logical contextuality rungs}

\begin{proposition}[Logical contextuality]\label{thm:logical-contextuality}\lean{Obsmetrologia.QuantumContextuality.LogicallyContextual}\leanok\uses{thm:possibilistic-model}
Logical contextuality asks whether at least one supported local section fails
to extend to any global section.
\end{proposition}

\begin{theorem}[The PR box is logically contextual]\label{thm:pr-logical-contextuality}\lean{Obsmetrologia.QuantumContextuality.prBox_logicallyContextual}\leanok\uses{thm:logical-contextuality, thm:pr-box-strong-contextual}
Strong contextuality of the PR box inhabits the logical contextuality rung.
\end{theorem}

\begin{theorem}[Hardy inhabits the strict middle rung]\label{thm:hardy-logical-not-strong}\lean{Obsmetrologia.QuantumContextuality.hardyBox_logical_not_strong}\leanok\uses{thm:logical-contextuality}
The Hardy support model has a global section and is not strongly contextual,
but a supported local section still extends to no global section.
\end{theorem}

\begin{theorem}[Hardy has a no-signaling probabilistic support]\label{thm:hardy-probabilistic-logical-not-strong}\lean{Obsmetrologia.QuantumContextuality.hardyBoxP_support_logical_not_strong}\leanok\uses{thm:hardy-logical-not-strong, thm:global-dist-no-signaling}
The Hardy phenomenon lifts to a no-signaling weighted empirical model whose
support is logically contextual but not strongly contextual.
\end{theorem}
